\documentclass[../../main.tex]{subfiles}
\begin{document}

\begin{flushright}
  \footnotesize\textit{【自動文字起こし・要確認】}
\end{flushright}

% Source: handwritten.pdf 第1問．
{\bf [解]}
\subparagraph{(1)}
\begin{align}
  \int_0^2\dfrac{2x}{\sqrt{x^2+4}}\,dx
\end{align}
において，$t=\sqrt{x^2+4}$とすると，
\begin{align}
  \dfrac{dt}{dx}=\dfrac{x}{t},\qquad
  \begin{array}{c|cc}x&0&2\\\hline t&2&2\sqrt{2}
  \end{array}
\end{align}
から
\begin{align}
  \int_0^2\dfrac{2x}{\sqrt{x^2+4}}\,dx
  =\int_2^{2\sqrt{2}}2\,dt
  =2[t]_2^{2\sqrt{2}}=4(\sqrt{2}-1).\label{eq:1}
\end{align}
又，
\begin{align}
  \begin{split}
    \int_0^2\dfrac{1}{\sqrt{x^2+4}}\,dx
    &=\int_0^{\pi/4}\dfrac{\dfrac{2}{\cos^2\theta}}{\dfrac{2}{\cos\theta}}\,d\theta
    \qquad(x=2\tan\theta)\\
    &=\int_0^{\pi/4}\dfrac{1}{\cos\theta}\,d\theta
    =\int_0^{\pi/4}\dfrac{\cos\theta}{1-\sin^2\theta}\,d\theta\\
    &=\int_0^{\pi/4}\left(\dfrac{1}{1-\sin\theta}+\dfrac{1}{1+\sin\theta}\right)
    \dfrac{\cos\theta}{2}\,d\theta\\
    &=\dfrac{1}{2}\left[-\log(1-\sin\theta)+\log(1+\sin\theta)\right]_0^{\pi/4}\\
    &=\dfrac{1}{2}\left\{-\log\left(1-\dfrac{\sqrt{2}}{2}\right)
    +\log\left(1+\dfrac{\sqrt{2}}{2}\right)\right\}.
  \end{split}\label{eq:2}
\end{align}
\cref{eq:1,eq:2}から，
\begin{align}
  (\text{与式})=4(\sqrt{2}-1)+\dfrac{1}{2}\log\dfrac{2+\sqrt{2}}{2-\sqrt{2}}.
\end{align}

\subparagraph{(2) 第1の方法}
$n$段の昇り方のうち，最後で1段昇ったものを$a_n$，2段昇ったものを$b_n$と表す．
\begin{figure}[H]
  \begin{center}
    \begin{tikzpicture}[>=Stealth,x=1.6cm,y=0.8cm]
      \node at (0,2) {$n$}; \node at (1,2) {$n+1$};
      \node (a) at (0,1) {$a_n$}; \node (b) at (0,0) {$b_n$};
      \node (c) at (1,1) {$a_{n+1}$};
      \draw[->] (a)--(c); \draw[->] (b)--(c);
      \node at (2.5,2) {$n$}; \node at (3.5,2) {$n+1$}; \node at (4.5,2) {$n+2$};
      \node (d) at (2.5,1) {$a_n$}; \node (e) at (4.5,0) {$b_{n+2}$};
      \draw[->] (d)--(e);
    \end{tikzpicture}
  \end{center}
  \caption{最後に昇る段数による昇り方の対応}
\end{figure}
上図から
\begin{align}
  \begin{cases}a_{n+1}=a_n+b_n,\\b_{n+2}=a_n.
  \end{cases}
\end{align}
これと$a_1=1$，$b_1=0$，$b_2=1$から，$a_n$，$b_n$を表出すれば
\begin{align}
  \begin{array}{c|rrrrrrrrrrrrrr}
    n&2&3&4&5&6&7&8&9&10&11&12&13&14&15\\\hline
    a_n&1&2&3&4&6&9&13&19&28&41&60&88&129&189\\
    b_n&1&1&1&2&3&4&6&9&13&19&28&41&60&88
  \end{array}
\end{align}
よって，求める場合の数は
\begin{align}189+88=277.
\end{align}

\subparagraph{第2の方法}
$A$，$B$を，以下の条件をみたすように同じ大きさのマスに並べればよい．
\begin{align}
  &a+2b=15\quad(a,b\text{は}A,B\text{の数}),\\
  &B\text{はとなりあわない}.
\end{align}
$(a,b)=(15,0),(13,1),(11,2),(9,3),(7,4),(5,5)$である．
それぞれの場合の数は
\begin{align}
  1\text{通り},\quad14\text{通り},\quad
  {}_{12}C_2=66\text{通り},\quad{}_{10}C_3=120\text{通り},\quad
  {}_8C_4=70\text{通り},\quad{}_6C_5=6\text{通り}.
\end{align}
よって
\begin{align}1+14+66+120+70+6=277.
\end{align}

\end{document}
